\begin{theorem}[Regular torus parent constraints force plaquette flatness]%
    \label{thm:peps_regular_torus_parent_flatness}
    \lean{TNLean.PEPS.IsGInjective.torusPlaquetteHolonomy_eq_one_of_parent_annihilates,
        TNLean.PEPS.IsGInjective.regionParent_not_annihilates_torusPlaquetteFlux,
        TNLean.PEPS.IsGInjective.torusPlaquetteHolonomy_eq_one_of_parentHamiltonian_annihilates,
        TNLean.PEPS.torusPlaquetteHolonomy_eq_one_iff_squareTransport,
        TNLean.PEPS.IsGInjective.isTorusFlat_of_parentHamiltonian_annihilates_twistedState}
    \leanok
    \uses{thm:peps_regular_region_parent_flatness,
        def:peps_region_parent_hamiltonian}
    Let a native four-leg regular $G$-injective tensor be placed on
    a square torus with both periods at least three. Insert arbitrary
    regular group operators on its ordered bonds, and suppose that
    the resulting physical PEPS vector is nonzero. If an original
    plaquette parent term annihilates this vector, the holonomy
    around that plaquette is the identity.

    If the sum of all original plaquette parent terms annihilates
    it, every elementary plaquette has identity holonomy. Equivalently,
    writing $A_v$ for transport to the right and $B_v$ for transport
    upward, one has
    \[
        B_{v+(1,0)}A_v=A_{v+(0,1)}B_v
    \]
    at every vertex. In particular, a nonzero one-bond flux state
    with nonidentity flux is not annihilated by the original parent
    term surrounding its endpoint.

    No $G$-isometry is assumed. This is a necessary local step in
    \cite[Theorem~5.5]{Schuch2010PEPS}; it does not claim that
    arbitrary parent-kernel vectors have inserted-bond descriptions,
    or that flatness alone gives the full torus ground-space classification.
\end{theorem}

\begin{proof}\leanok
    Each native plaquette is the support of a four-step closed walk,
    so its induced graph is connected. Apply local regular parent
    flatness to the induced walk and then identify its holonomy with
    the original native plaquette holonomy. Positivity of every
    regional parent term implies that a vector killed by their sum
    is killed by each term separately.

    The ordered plaquette holonomy is
    $B_v^{-1}A_{v+(0,1)}^{-1}B_{v+(1,0)}A_v$.
    It is the identity exactly when the two elementary square
    transports agree. The explicit one-bond insertion has its
    prescribed group element as holonomy, giving the exclusion of
    nonidentity flux.
\end{proof}
